\section{L'autore, il vuoto normativo e le stime dei danni}\label{sec:Autore}
In questa sezione si analizzano l'autore del worm, le conseguenze legali e le stime dei danni.

\subsection{Onel de Guzman}\label{subsec:deGuzman}
Le indagini portarono in pochi giorni a Manila, dove i provider riuscirono a risalire alla linea telefonica da cui era partito il worm.
L'utenza apparteneva all'appartamento di Onel de Guzman, uno studente di 23 anni dell'AMA Computer College, che venne arrestato l'8 maggio 2000~\cite{voa2000}.
Durante la perquisizione furono trovati diciassette dischetti, e uno di questi conteneva un programma molto simile al worm~\cite{ecommercetimes2000a}.

Un altro indizio importante arrivò proprio dall'università.
Infatti qualche mese prima de Guzman aveva proposto come tesi un programma per rubare account e password di accesso a Internet, presentandolo come un modo per permettere a tutti di navigare senza pagare.
La proposta venne respinta, e secondo il college anche due sue bozze precedenti riguardavano altre parti del virus~\cite{outlaw2000}.
Furono quindi i docenti a notare la somiglianza tra i progetti rifiutati e le funzioni di ILOVEYOU~\cite{computable2000}.

Nonostante le prove contro di lui, de Guzman si proclamò sempre innocente, arrivando solo a parlare della possibilità che il virus si fosse diffuso per sbaglio~\cite{outlaw2000}.
Solo nel 2020, vent'anni dopo il fatto, ha ammesso di esserne l'autore~\cite{mimms,bbc2020}.

\subsection{Il vuoto normativo}\label{subsec:vuotoNormativo}
Nel diritto filippino, come in quello italiano, vale il principio di legalità, ovvero non esiste reato senza una legge che lo preveda~\cite[pp.~34--35]{malibiran}, e nel maggio 2000 le Filippine non avevano leggi sui crimini informatici~\cite{deseret2000}.
Gli investigatori, notata la lacuna, ripiegarono su reati pensati per tutt'altro: furto e violazione di una legge sulle frodi con carte di credito~\cite{voa2000,seclists2000}.

Per rimediare, il 14 giugno 2000 il Paese approvò in fretta l'\emph{Electronic Commerce Act} (Republic Act 8792), che punisce l'accesso non autorizzato ai sistemi informatici e la diffusione di virus~\cite[p.~651]{sy2015}.
La legge però non era retroattiva e non poté essere usata contro de Guzman~\cite{deseret2000}.
Caddero poi anche le accuse di ripiego, perché la legge sulle carte di credito non riguardava l'hacking e le prove per il furto non bastavano~\cite{deseret2000,computerworld2000}.
De Guzman, quindi, non venne mai processato.

\subsection{Le stime dei danni}\label{subsec:stimeDanni}
De Guzman non ricavò nulla dall'attacco: lo scopo del worm era raccogliere password per navigare gratis, e gli account che dovevano riceverle vennero disattivati già nelle prime ore~\cite[p.~2]{gao2000iloveyou}.
Per le vittime, al contrario, le stime più citate parlano di danni per miliardi di dollari.

Fin da subito ci furono difficoltà nello stimare correttamente i danni causati dal worm.
Pochi giorni dopo l'attacco, il General Accounting Office (GAO) riportò al Senato statunitense stime comprese tra 100 milioni e oltre 10 miliardi di dollari~\cite[pp.~4--5]{gao2000iloveyou}, spiegando che molte perdite sono ardue da misurare e che le aziende tendono a nasconderle~\cite[p.~5]{gao2000iloveyou}.

La cifra più citata, circa 8,7 miliardi, fu fornita dalla società di consulenza Computer Economics.
Comprendeva oltre 50 varianti del worm e sommava al costo di pulizia il tempo di lavoro perso, che è la voce più incerta~\cite{outlaw2001}.

\subsection{Considerazioni}
I danni documentati furono reali, come i 1.000 file danneggiati alla NASA o le 2.800 ore di lavoro per il ripristino al Dipartimento del Lavoro~\cite[pp.~7--8]{gao2000iloveyou}.
Si misurano però in ore e file, non in miliardi, quindi si deve considerare la cifra in dollari come un ordine di grandezza.

Il vuoto normativo e le stime così incerte mostrano la stessa cosa: nel 2000 né le leggi né le istituzioni erano pronte ad affrontare un attacco informatico su scala mondiale.
